\documentclass[12pt]{article}
\usepackage{notes-common}
\usepackage{graphicx}
\geometry{margin=0.35in}
\AtBeginDocument{\fontsize{14}{17}\selectfont}

% Two short worked examples (Knight 29.3/29.5 territory): finding the field
% in a velocity selector (solving the force-balance condition for B instead
% of v), and the energy difference between two orientations of a magnetic
% dipole in a uniform field. See
% velocity-selector-dipole-energy-notes-key.tex for the filled-in version.
\newcommand{\eqline}[2][-0.5mm]{%
  \par\vspace{1mm}\noindent #2\par\vspace{#1}%
  \noindent\textcolor{rulehue}{\rule{\linewidth}{0.5pt}}\par\vspace{4mm}%
}

\begin{document}

\noteshead{Magnetic Field from a Force Balance, and Dipole Energy}{}

\goal{
\begin{itemize}[itemsep=5pt,parsep=2pt,topsep=2pt]
\item Apply the velocity-selector force balance to solve for the
\emph{field} instead of the speed --- same equation, different unknown.
\item Use $U=-\vec\mu\cdot\vec B=-\mu B\cos\theta$ to find a dipole's
potential energy at a given orientation, and the \emph{change} in
energy between two different orientations.
\end{itemize}}

%============================================================
\nsec{Worked example: finding the field in a velocity selector}

Recall the velocity selector: a charge moving through crossed
$\vec E$ and $\vec B$ fields goes straight only when the electric and
magnetic forces exactly cancel, $qE=qvB$. Last time you solved that
for the speed $v$. This time you're given the speed and have to solve
for the field.

\textit{A charged particle moves with speed $v=3.0\times10^5$ m/s
through a region containing a uniform electric field
$E=6.0\times10^4$ N/C, directed vertically downward. The particle
passes straight through without deflecting. What is the magnitude of
the magnetic field $\vec B$ present in the region? Assume $\vec B$ is
perpendicular to both $\vec v$ and $\vec E$, and the particle is
positively charged.}

\textbf{MODEL.} The particle goes straight only if the electric force
and the magnetic force on it are \bl[3cm], exactly canceling.

\textbf{SOLVE.} Set the force magnitudes equal:

\[
qE = \bl[3cm]
\]

The charge $q$ \bl[3cm] --- the balance condition doesn't depend on
it. Solve for $B$:

\result[1.6cm]{$B = $}

Plug in the numbers:

\eqline[0.2cm]{$B = $}

\textbf{Direction.} The electric force on the positive charge points
downward (along $\vec E$), so the magnetic force $q\vec v\times\vec B$
must point \bl[2.5cm] to cancel it. With $\vec v$ horizontal, the
right-hand rule then fixes which way $\vec B$ must point.

\checkpoint{If the particle's speed were doubled but $E$ stayed the
same, would the required $B$ go up, go down, or stay the same? Does
that match your intuition --- does a faster particle need a
\emph{stronger} or \emph{weaker} field to still be balanced?}

%============================================================
\nsec{Worked example: energy of a magnetic dipole, two orientations}

Recall the energy of a magnetic dipole in a field: $U=-\vec\mu\cdot
\vec B=-\mu B\cos\theta$, lowest when $\vec\mu$ is aligned with
$\vec B$ ($\theta=0$) and highest when anti-aligned ($\theta=180^\circ$).

\textit{A magnetic dipole with moment $\mu=0.60\text{ A}{\cdot}\text{m}^2$
sits in a uniform field $B=0.30$ T. Someone rotates it from
$\theta_1=40^\circ$ to $\theta_2=110^\circ$ (both measured from
$\vec B$). What is the \emph{change} in the dipole's potential energy?}

\textbf{MODEL.} The energy change is just the energy in the final
orientation minus the energy in the initial one:
$\Delta U = \bl[3.5cm]$.

\textbf{SOLVE.} Find the energy at each angle separately:

\[
U_1 = -\mu B\cos\theta_1 = \bl[6cm]
\]

\[
U_2 = -\mu B\cos\theta_2 = \bl[6cm]
\]

Now subtract:

\result[1.8cm]{$\Delta U = U_2-U_1 = $}

Plug in the numbers:

\eqline[0.2cm]{$\Delta U = $}

\checkpoint{Is $\Delta U$ positive or negative here? $\theta_1=40^\circ$
is closer to alignment ($0^\circ$) than $\theta_2=110^\circ$ is. Does
the sign of $\Delta U$ match the fact that torque always twists a
dipole \emph{toward} alignment --- so moving \emph{away} from alignment
should take energy \emph{in}, not give it up?}

\end{document}
